\section{模的直和}

记号约定:$A$是环,$\left(E_{i}\right)_{i\in I},\left(F_{i}\right)_{i\in I},\left(G_{i}\right)_{i\in I}$是$A$-模构成的族.

\begin{definition}{模的外直和}{}
	定义$\left(M_{i}\right)_{i\in I}$的外直和为$\bigboxplus\limits_{i\in I}M_{i}\coloneq\left\{\left(x_{i}\right)_{i\in I}\in\prod\limits_{i\in I}E_{i}\middle|\left(x_{i}\right)_{i\in I}\text{有限支撑}\right\}$.
\end{definition}

\begin{remark}{记号说明}{}
	\begin{enumerate}
		\item $\bigboxplus\limits_{i\in I}E_{i}$是$\prod\limits_{i\in I}E_{i}$的子模,当$I$有限时$\bigboxplus\limits_{i\in I}E_{i}=\prod\limits_{i\in I}E_{i}$.
		\item 若$I=\llbracket p,q\rrbracket$是$\Z$的区间,则记$E_{p}\boxplus E_{p+1}\boxplus\ldots\boxplus E_{q}\coloneq\bigboxplus\limits_{i\in I}E_{i}$.
	\end{enumerate}
\end{remark}

\begin{definition}{外直和的典范嵌入}{}
	对诸$j\in J$,定义$\iota_{j}:M_{j}\hookrightarrow\bigboxplus\limits_{i\in I}M_{i},x_{j}\mapsto\left(x_{i}\right)_{i\in I},x_{i}=
	\begin{cases}
		x_{j},&i=j,\\
		0,&i\neq j.
	\end{cases}$称$\left(\iota_{i}\right)_{i\in I}$是$\bigboxplus\limits_{i\in I}M_{i}$的典范嵌入族.
\end{definition}

\begin{proposition}{典范嵌入和典范满射之间的关系}{}
	\begin{enumerate}
		\item $\left(\iota_{i}\right)_{i\in I}$是一族$A$-线性映射.
		\item 对任意$i,j\in I,\pr_{i}\circ\iota_{j}=
		\begin{cases}
			\id_{E_{i}},&i=j,\\
			0,&i\neq j.
		\end{cases}$
		\item 对任一$i\in I$,$\im\left(\iota_{i}\right)\simeq E_{i}$.我们称$\im\left(\iota_{i}\right)$为$E$的第$i$个组件.
		\item 对任一$x\in\bigboxplus\limits_{i\in I}E_{i}$,有$\sum\limits_{i\in I}\iota_{i}\left(\pr_{i}\left(x\right)\right)=x$.
	\end{enumerate}
\end{proposition}

\begin{proposition}{直和的泛性质}{universal-property-of-direct-sum}
	设$\forall i\in I\left(f_{i}\in\Hom_{A}\left(E_{i},F\right)\right)$,则对任一$i\in I$下图交换.
	\begin{center}
		\begin{tikzcd}
			{E_{i}} && {\bigboxplus\limits_{i\in I}E_{i}} \\
			& F
			\arrow["{\iota_{i}}", hook, from=1-1, to=1-3]
			\arrow["f_{i}"', from=1-1, to=2-2]
			\arrow["f", dashed, from=1-3, to=2-2]
			\end{tikzcd}
	\end{center}
	具体的,$f:\bigboxplus\limits_{i\in I}E_{i}\to F,x\mapsto\sum\limits_{i\in I}f_{i}\left(\pr_{i}\left(x\right)\right)$.
\end{proposition}

\begin{definition}{线性映射族的和}{}
	设$\forall i\in I\left(f_{i}\in\Hom_{A}\left(E_{i},F\right)\right)$.称$\sum\limits_{i\in I}f_{i}:\bigboxplus\limits_{i\in I}E_{i}\to F,\left(x_{i}\right)_{i\in I}\mapsto\sum\limits_{i\in I}x_{i}$为$\left(f_{i}\right)_{i\in I}$的和.
\end{definition}

\begin{proposition}{}{}
	设$J\subseteq I$,则典范嵌入$\iota_{J}:\bigboxplus\limits_{i\in J}E_{i}\hookrightarrow\bigboxplus\limits_{i\in I}E_{i}$是$A$-线性映射.
\end{proposition}

\begin{proposition}{直和的``结合律''}{}
	设$I\neq\emptyset$且有划分$\left(I_{\lambda}\right)_{\lambda\in\Lambda}$,则有典范同构$\bigboxplus\limits_{i\in I}E_{i}\simeq\bigboxplus\limits_{\lambda\in\Lambda}\bigboxplus\limits_{i\in I_{\lambda}}M_{i}$.
\end{proposition}

\begin{proposition}{直和与直积的$\Hom$函子的性质}{}
	设$\left(F_{j}^{\prime}\right)_{j\in J}$是一族$A$-模,则$\varphi:\Hom_{A}\left(\bigboxplus\limits_{i\in I}E_{i},\prod\limits_{j\in J}F_{j}^{\prime}\right)\to\prod\limits_{\left(i,j\right)\in I\times J}\Hom_{A}\left(E_{i},F_{j}^{\prime}\right),g\mapsto\left(\pr_{j}\circ f\circ\iota_{i}\right)_{i\in I,j\in J}$是$\Z$-模同构,其中$\left(\pr_{j}\right)_{j\in J}$是$\prod\limits_{j\in J}F_{j}^{\prime}$的典范投影族,$\left(\iota_{i}\right)_{i\in I}$是$\bigboxplus\limits_{i\in I}E_{i}$的典范嵌入族.
\end{proposition}

\begin{proposition}{直和分解存在的等价条件}{}
	设$\forall i\in I\left(f_{i}\in\Hom_{A}\left(E_{i},F\right)\right)$,$E=\bigoplus\limits_{i\in I}E_{i},f=\sum\limits_{i\in I}f_{i}$.下列陈述等价:
	\begin{enumerate}
		\item $f\in\Isom_{A}\left(E,F\right)$.
		\item 对任一$i\in I$,存在$g_{i}\in\Hom_{A}\left(F,E_{i}\right)$满足下列条件:
		\begin{itemize}
			\item 对任意$i,j\in I,g_{i}\circ f_{j}=
			\begin{cases}
				\id_{E_{i}},&j=i,\\
				0,&j\neq i.
			\end{cases}$
			\item 对任意$y\in F$,$\left(g_{i}\left(y\right)\right)_{i\in I}$有限支撑且$\sum\limits_{i\in I}f_{i}\left(g_{i}\left(y\right)\right)=y$.
		\end{itemize}
	\end{enumerate}
\end{proposition}

\begin{proposition}{直和的函子性}{}
	设$E=\bigboxplus\limits_{i\in I}E_{i},F=\bigboxplus\limits_{i\in I}F_{i},\forall i\in I\left(f_{i}\in\Hom_{A}\left(E_{i},F_{i}\right)\right)$,则对任一$i\in I$下图交换(其中$\left(\iota_{i}^{\prime}\right)_{i\in I}$是$F$的典范嵌入族).
	\begin{center}
		\begin{tikzcd}
			{E_{i}} & E \\
			{F_{i}} & F
			\arrow["{\iota_{i}}", hook, from=1-1, to=1-2]
			\arrow["{f_{i}}"', from=1-1, to=2-1]
			\arrow["{\exists!f}", dashed, from=1-2, to=2-2]
			\arrow["{\iota_{i}^{\prime}}"', hook, from=2-1, to=2-2]
		\end{tikzcd}
	\end{center}
	具体的,$f:E\to F,\left(x_{i}\right)_{i\in I}\mapsto\left(f_{i}\left(x_{i}\right)\right)_{i\in I}$
\end{proposition}

\begin{proposition}{线性映射族的直和是线性映射}{}
	设$\forall i\in I\left(f_{i}\in\Hom_{A}\left(E_{i},F_{i}\right)\right)$,则是$A$-线性映射.
\end{proposition}

\begin{definition}{线性映射族的直和}{}
	设$\forall i\in I\left(f_{i}\in\Hom_{A}\left(E_{i},F_{i}\right)\right)$.称$\bigboxplus\limits_{i\in I}f_{i}:\bigboxplus\limits_{i\in I}E_{i}\to\bigboxplus\limits_{i\in I}M_{i},\left(x_{i}\right)_{i\in I}\mapsto \left(f_{i}\left(x_{i}\right)\right)_{i\in I}$为$\left(f_{i}\right)_{i\in I}$的乘积.
\end{definition}

\begin{proposition}{线性映射族的直和与正合性}{}
	设$\forall i\in I\left(f_{i}\in\Hom_{A}\left(E_{i},F_{i}\right)\wedge g_{i}\in\Hom_{A}\left(F_{i},G_{i}\right)\right)$.记$E=\bigboxplus\limits_{i\in I}E_{i},F=\bigboxplus\limits_{i\in I}F_{i},G=\bigboxplus\limits_{i\in I}G_{i},f=\bigboxplus\limits_{i\in I}f_{i},g=\bigboxplus\limits_{i\in I}g_{i}$,则
	\begin{center}
		$E\xlongrightarrow{f}F\xlongrightarrow{g}G$正合$\iff$对任一$i\in I$,$E_{i}\xlongrightarrow{f_{i}}F_{i}\xlongrightarrow{g_{i}}G_{i}$正合.
	\end{center}
\end{proposition}

\begin{corollary}{}{}
	设$\forall i\in I\left(f_{i}\in\Hom_{A}\left(E_{i},F_{i}\right)\right)$.记$E=\bigboxplus\limits_{i\in I}E_{i},F=\bigboxplus\limits_{i\in I}F_{i}$.
	\begin{enumerate}
		\item $\ker\left(f\right)=\bigboxplus\limits_{i\in I}\ker\left(f_{i}\right)$,$\im\left(f\right)=\bigboxplus\limits_{i\in I}\im\left(f_{i}\right)$.
		\item 存在典范同构$\coim\left(f\right)\simeq\bigboxplus\limits_{i\in I}\coim\left(f_{i}\right)$,$\coker\left(f\right)\simeq\bigboxplus\limits_{i\in I}\coker\left(f_{i}\right)$.
		\item $f$是单射(相对地,满射,双射)$\iff$ $\left(f_{i}\right)_{i\in I}$是一族单射(相对地,满射,双射).
	\end{enumerate}
\end{corollary}

\begin{proposition}{}{}
	设对任意$i\in I$,$F_{i}$是$E_{i}$的子模.记$E=\bigboxplus\limits_{i\in I}E_{i},F=\bigboxplus\limits_{i\in I}F_{i}$,则$F$是$E$的子模,并且有典范同构$\bigboxplus\limits_{i\in I}\left(E_{i}/F_{i}\right)\simeq E/F$.
\end{proposition}

\begin{corollary}{}{}
	设$\left(E_{i}^{\prime}\right)_{i\in I},\left(\tilde{F}_{j}\right)_{j\in J},\left(\tilde{F}_{j}^{\prime}\right)_{j\in J}$是三族$A$-模,$E=\bigboxplus\limits_{i\in I}E_{i},E^{\prime}=\bigboxplus\limits_{i\in I}E_{i}^{\prime},\tilde{F}=\prod\limits_{j\in J}\tilde{F}_{j},\tilde{F}^{\prime}=\prod\limits_{j\in J}\tilde{F}_{j}^{\prime}$,并且
	\begin{align*}
		\forall i\in I\left(u_{i}\in\Hom_{A}\left(E_{i},\tilde{F}_{i}\right)\right)\forall j\in J\left(v_{j}\in\Hom_{A}\left(E_{j}^{\prime},\tilde{F}_{j}^{\prime}\right)\right),
	\end{align*}
	则下图交换(其中$\phi,\phi^{\prime}$是典范同构).
	\begin{center}
		\begin{tikzcd}
			{\Hom_{A}\left(E^{\prime},\tilde{F}^{\prime}\right)} && {\prod\limits_{\left(i,j\right)\in I\times J}\Hom_{A}\left(E_{i}^{\prime},\tilde{F}_{j}^{\prime}\right)} \\
			\\
			{\Hom_{A}\left(E,\tilde{F}\right)} && {\prod\limits_{\left(i,j\right)\in I\times J}\Hom_{A}\left(E_{i},\tilde{F}_{j}\right)}
			\arrow["{\phi^{\prime}}", from=1-1, to=1-3]
			\arrow["{\prod\limits_{\left(i,j\right)\in I\times J}\Hom_{A}\left(\bigboxplus\limits_{i\in I}u_{i},\prod\limits_{j\in J}v_{j}\right)}", from=3-1, to=1-1]
			\arrow["\phi"', from=3-1, to=3-3]
			\arrow["{\prod\limits_{\left(i,j\right)\in I\times J}\Hom_{A}\left(u_{i},v_{j}\right)}"', from=3-3, to=1-3]
		\end{tikzcd}
	\end{center}
\end{corollary}

\begin{definition}{$I$-次直和}{}
	我们称$E^{\left(I\right)}\coloneq\bigboxplus\limits_{i\in I}E$为$E$的$I$-次直和.
\end{definition}

\begin{definition}{余对角映射}{}
	我们称$\nabla:E^{\left(I\right)}\to E,\left(x_{i}\right)_{i\in I}\mapsto\sum\limits_{i\in I}x_{i}$为余对角映射.
\end{definition}

\begin{remark}{}{}
	本节的所有内容都可以推广到不一定交换的群,只需要做如下替换:
	\begin{itemize}
		\item 模$\mapsto$群(乘法记号),直和$\mapsto$限制和,子模$\mapsto$正规子群,线性映射$\mapsto$群同态.
		\item 命题(\ref{prop:universal-property-of-direct-sum})增加下列条件:对任一$i,j\in I$且$i\neq j,x_{i}\in E_{i},x_{j}\in E_{j},f_{i}\left(x_{i}\right)$和$f_{j}\left(x_{j}\right)$可交换.
	\end{itemize}
\end{remark}




















